%%%%%%%%%%%%%%%%%%%%%%%%%%%%%%%%%%%%%%%%%%%%%%%%%%%%%%%%%%%%%%%%%%%%%%%%%%%%%%%%
%%
\cleardoubleoddpage%  Make sure to start each chapter on a new odd page
\chapter{Data Format}
\label{sec:data-format}
\section{Word definitions}
\subsection*{What is a "task"}
Before we can specify how we define a task, we first need to declare what a task actually is. We define a task as single event in time, with a set in duration, non reoccurring and mostly, independent. We say "mostly" here, as a task can be time dependent, on another task.

Furthermore, we decided for these restrictions to reduce our set of movable parts in the equations to reduce computational intensity later in the process, while also ensuring ease of use for defining a task.

As defining tasks unambiguously with just a set of constraints is almost impossible we restrict ourselves to handle only static tasks with static time definitions. Meaning, a task has a set duration, and defines different time slots where such a task can be scheduled. This time slots - this is the static part - is defined using a 0 based time slot "array" in form of intervals. It starts counting from the 0 (now) up to a given horizon and every task inside that needs to be scheduled has defined possible time slots (intervals), which are also static - relative to the given beginning of the tasks (time slot 0).

\subsection*{What is a "time slot"}
Instead of restricting ourselves on minutes or equivalent time units, we define time slots to be an any unit of time. This allows us to solve arbitrarily granular problems.

\section{Input fields}
In this section we list our final input fields with our motivation why we need it and why it is better than other options we accounted for.
\subsection*{Horizon}
This \textit{horizon} field is solely to define the length of time the solver can define tasks to - the amount of time slots.

\subsection*{Tasks}
For convenience reasons - accessing the \textit{tasks} later on - and because a task needs to be uniquely identifiable, tasks is a dictionary with a unique string key and the following fields as a task object.
\subsubsection*{Duration}
Duration defines the length of a task in form of number of time slots as an integer.

\subsubsection*{Importance}
\label{subsubsec:data-format:importance}
We define \textit{importance} to be an integer (or infinity), which defines how important a task is. Importance defines the cost of an unscheduled task. A task that must be scheduled without exception should be assigned an importance value of or approaching infinity. In contrast, a task whose inclusion is not critical can be given a minimal importance score, close to 0.

\subsubsection{Time slots}
\label{subsubsec:data-format:time-slots}
Describing when a given task can be scheduled is represented by the \textit{timeSlots} field. Without it, one could not define when to, when to best or when not to schedule a task. We came up with the following options on how to efficiently define the field:
\begin{itemize}
    \item array of bits,
    \item array of interval
    \item and array of interval-cost pairs.
\end{itemize}
Bit-arrays are very easy to handle and efficient to store - but as storage is not a concerning factor, and the representation scales linearly with the horizon it becomes harder to understand for humans, the bigger the horizon. Therefore, we use array of intervals as it can be easily converted to an array of bits\footnote{Interval $[x,y] = [b_0, b_1, \dots]$ where $b_i = \begin{cases}0, & 0 \leq i < x\\ 1 & x \leq i \leq y\\ 0 & \text{else}\end{cases}$ as a bit array}, while still being easily human-readable. For expressiveness, we additionally added a cost for each interval and consequently stick with an array of interval-cost pairs.

\subsection*{Dependencies}
In the \textit{dependencies} we depict preceding tasks. It in tern contains again intervals in the same format as in \nameref{subsubsec:data-format:time-slots}. Additionally, similarly to the \nameref{subsubsec:data-format:importance} of a task, we define \textit{unmetCost} if a task does not conform to its dependencies. The necessity of the here described fields are already described in \nameref{subsubsec:data-format:time-slots} and \nameref{subsubsec:data-format:importance}.

\subsubsection*{Location}
Planning a timetable well improves throughput while simultaneously reducing travel time by grouping tasks in the same area together. This is what \textit{location} is used for. For this we define the location by a unique integer alongside with \textit{unmetCost}, for when the grouping is not well-managed. This is the most basic set of fields, one could use for this purpose. For the general functionality it is not required, but again improves expressively by a lot, while maintaining a minimalistic language.

\section{JSON Definition}
For ease of use we decided against defining our own language. Therefore, we use \ac{JSON} to represent our basic language. The decision fell on this format, as it is widely spread and easily human-readable. If one wanted to instead use XML, it can be converted to JSON and back \cite{bib:2017-LaFontaine}.
\begin{lstlisting}[language=json]
{
    "horizon": 100,
    "tasks": {
        "SomeTaskId": {
            "name": "Some Cleartext name",
            "duration": 5,
            "importance": 100,
            \textcolor{red}{TODO}"urgency": 5,
            "timeSlots": [
                { "interval": [5,12], "cost": 5},
                { "interval": [20,22], "cost": 0},
                { "interval": [20,28], "cost": 25},
            ],
            "dependencies": [
                { 
                    "task": "AnotherTaskId",
                    "intervals": [
                        { "interval": [10, 20], "cost": 5},
                        { "interval": [20, 1110], "cost": 15}
                    ],
                    "unmetCost": 100
                }
            ],
            "location": { "id": 1, "unmetCost": 150 }
        }
    }
}
\end{lstlisting}

%% 
%%%%%%%%%%%%%%%%%%%%%%%%%%%%%%%%%%%%%%%%%%%%%%%%%%%%%%%%%%%%%%%%%%%%%%%%%%%%%%%%
